\section{Conclusiones}
\label{sec:conclusiones}

PARTS\_SORTING\_PLATFORM demuestra que un mismo problema de
identificación de piezas puede resolverse con dos arquitecturas
distintas según el entorno de despliegue disponible: una orientada a
múltiples usuarios sobre infraestructura administrada en la nube
(CLOUD\_MODE), y otra orientada a un único banco de pruebas físico sin
acceso a internet (LOCAL\_MODE). La separación de responsabilidades
entre el control físico de la banda (habilitado únicamente mediante
huella digital autorizada), la detección automática de la pieza bajo
la cámara, y el pipeline de clasificación y medición de diámetro,
permite que cada componente evolucione de forma independiente. El
modelo YOLOv8 entrenado específicamente para el banco de pruebas local
constituye la vía principal de clasificación en LOCAL\_MODE, mientras
que CLOUD\_MODE aún opera sobre un clasificador de referencia pendiente
de calibración. Queda como trabajo futuro calibrar los perfiles de
referencia en CLOUD\_MODE y ampliar el conjunto de entrenamiento del
modelo local para las clases con menor representación.
